% !TEX program = xelatex
% Compile this file with XeLaTeX to use Times New Roman via fontspec.
\documentclass[10pt,letterpaper]{article}

\usepackage[margin=0.5in]{geometry}
\usepackage{fontspec}
\usepackage{xcolor}
\definecolor{linkblue}{RGB}{0,86,179}
\usepackage{enumitem}
\usepackage[colorlinks=true, urlcolor=linkblue, linkcolor=linkblue]{hyperref}
\usepackage{titlesec}
\usepackage{setspace}
\usepackage[normalem]{ulem}

\setmainfont{Times New Roman}

\pagestyle{empty}
\setlength{\parindent}{0pt}
\setlength{\parskip}{0pt}
\setlength{\tabcolsep}{0pt}
\setstretch{0.93}
\raggedbottom

\titleformat{\section}{\large\bfseries}{}{0em}{}[\titlerule]
\titlespacing{\section}{0pt}{4pt}{2pt}

\setlist[itemize,1]{
    leftmargin=1.1em,
    label=•,
    topsep=0pt,
    itemsep=0pt,
    parsep=0pt,
    partopsep=0pt
}
\setlist[itemize,2]{
    leftmargin=1.55em,
    label=◦,
    topsep=0pt,
    itemsep=0pt,
    parsep=0pt,
    partopsep=0pt
}

\newcommand{\ressection}[1]{\section*{#1}}

\newcommand{\resheading}[4]{%
    \textbf{#1}\hfill #2\par
    \textit{#3}\hfill #4\par
}

\newcommand{\translabel}[1]{%
    {\setlength{\ULdepth}{2.2pt}\uline{#1}}%
}

\begin{document}
\fontsize{10.5pt}{11.8pt}\selectfont

\begin{center}
    {\LARGE \textbf{Yincheng Zhou}} \\[2pt]
    (267) 788-2488 | \href{mailto:yzhou29@seas.upenn.edu}{yzhou29@seas.upenn.edu} | \href{https://github.com/ArtysicistZ}{github.com/ArtysicistZ} | \href{https://artysicistz.github.io}{artysicistz.github.io}
\end{center}
\vspace{-8pt}

\ressection{EDUCATION}

\resheading{University of Pennsylvania}{Philadelphia, PA}{Vagelos Integrated Program in Energy Research (VIPER)}{September 2025 -- Present}
\begin{itemize}
    \item \textbf{Dual Degree}: B.S.E. in Computer Science (AI Concentration) | B.A. in Physics \& Astronomy
    \item \textbf{Relevant Coursework:} Big Data Analysis, Computer Systems, Data Structures \& Algorithms, OOP
    \item \textbf{GPA:} 4.0 / 4.0
\end{itemize}

\ressection{EXPERIENCE}

\textbf{GRASP Laboratory} \;|\; \textit{University of Pennsylvania}\hfill Philadelphia, PA \par
\textit{Machine Learning Researcher, Advisor: Prof. Pratik Chaudhari}\hfill January 2026 -- Present
\begin{itemize}
    \item Proposed SPD, a sparse monocular depth estimator that predicts depth only at queried pixels, achieving 6$\times$ inference speedup over dense methods while preserving accuracy (AbsRel 0.11 on NYU Depth V2) for real-time robotics.
    \item Developed a dense-train / sparse-infer paradigm with a 350K-param split decoder and cross-attention fusion on DINOv2 features, proving sparse output matches dense pixel-wise prediction through systematic ablation across 15 architecture variants.
\end{itemize}
\vspace{1pt}
\resheading{Welta, Inc. \;|\; \normalfont\textit{AI-powered cross-platform ad optimization engine} \textbf{\href{https://app.welta.ai}{[Demo]}}}{Remote}{Software Engineer Intern}{March 2026 -- Present}
\begin{itemize}
    \item Designed ad spend optimizer using Bayesian response curves and convex optimization to reallocate budget across channels at fixed cost; validated on 93 e-commerce brands, delivering 23\% higher return on ad spend comparing to human baseline.
    \item Engineered ad simulation engine that tokenizes creatives and models persona-level reactions via LLM; reduced inference latency to 20\% of baseline while improving prediction accuracy by 55\%.
\end{itemize}
\vspace{1pt}
\textbf{Human Machine Intelligence Laboratory} \;|\; \textit{Peking University}\hfill Remote \par
\textit{Machine Learning Researcher, Advisor: Prof. Shanghang Zhang}\hfill March 2026 -- Present
\begin{itemize}
    \item Architected GUI agent post-training pipeline that generates step-level SFT data with MCTS and difficulty-weighted loss, outperforming agentic RL (GRPO) at 25\% of compute; raised OSWorld success rate from 12.3\% to 17.7\% (+44\% relative).
    \item Engineered distributed rollout infrastructure across 112 Docker VMs and 8 A100 GPUs with fault-tolerant container monitoring and action-replay VM cloning; collected 2,927 trajectories and curated 35K steps via 200 parallel LLM auditing agents.
\end{itemize}
\vspace{1pt}
\resheading{Franklink, Inc. \;|\; \normalfont\textit{World's first AI-native networking platform} \textbf{\href{https://artysicistz.github.io/projects/franklink-demo.html}{[Demo]}}}{Philadelphia, PA}{Founding Software Engineer}{October 2025 -- February 2026}
\begin{itemize}
    \item Designed iMessage-based multi-agent orchestration layer with persisted memory, selective dispatch across 5 sub-agents, and agentic RAG over 8 permissioned tools (DB, email, calendar); reduced user-facing latency by 30\%.
    \item Engineered a Kafka-backed event pipeline with at-least-once delivery, idempotency keys, and retry + DLQ; sustained 50+ concurrent sessions with p99 latency within 40\% of single-session baseline; deployed on AWS ECS.
\end{itemize}

\ressection{PROJECTS}

\textbf{AlphaOne} \;|\; \textit{Subject-oriented financial sentiment platform} \textbf{\href{https://github.com/ArtysicistZ/AlphaOne.git}{[GitHub]} \href{https://artysicistz.github.io/projects/alphaone-demo.html}{[Demo]} \href{https://huggingface.co/ArtysicistZ/deberta-absa-v2}{[HuggingFace]}}\hfill August 2025 -- February 2026
\begin{itemize}
    \item Fine-tuned DeBERTa for per-ticker aspect-based sentiment in multi-entity financial text, achieving 82.5\% accuracy (+65pp over FinBERT); curated 6,287 training triples via Reddit scraping, LLM labeling, and manual audit of 4,501 pairs.
    \item Engineered a 3-pass batch NLP pipeline (split, infer, commit) with claim-based row locking for concurrent-safe ingestion across 7 subreddits; serves real-time per-ticker trends through a full-stack dashboard.
\end{itemize}
\vspace{1pt}
\textbf{F-Zero Racing Agent} \;|\; \textit{Reinforcement learning for competitive SNES racing} \textbf{\href{https://github.com/ArtysicistZ/FZero}{[GitHub]}}\hfill March 2026 -- April 2026
\begin{itemize}
    \item Trained a PPO agent across 80 parallel SNES emulators with custom reward shaping to race F-Zero, reaching global top 300 on a task previously deemed unsolvable due to sparse terminal-only rewards over 3,000+ step horizons.
    \item Engineered dense reward from raw RAM by spline-interpolating 58 checkpoints into a 660-point track centerline; designed dual-input CNN+MLP policy fusing 4-frame visual input with 59-dim state features.
\end{itemize}
\vspace{1pt}
\textbf{IncidentGuard} \;|\; \textit{Predict-then-act multi-agent incident remediation} \textbf{\href{https://github.com/hanshengzhu0001/stratusxpennai}{[GitHub]}}\hfill April 2026
\begin{itemize}
    \item Built a multi-agent system that predicts remediation outcomes before execution, rejecting dangerous reflexes via counterfactual ranking; reduced simulated recovery latency by 46\% over naive baseline across 4 failure scenarios.
    \item Implemented closed-loop plan-execute-verify pipeline with live Prometheus metrics, OpenClaw browser automation, and DuckDB case persistence; system gracefully degrades across 3 fallback paths when infrastructure is unavailable.
\end{itemize}
\vspace{1pt}
\textbf{NL2SQL Bot} \;|\; \textit{Natural language to SQL with agentic orchestration} \textbf{\href{https://github.com/ArtysicistZ/NL2SQL_Bot.git}{[GitHub]} \href{https://artysicistz.github.io/projects/nl2sql.html}{[Demo]}}\hfill December 2025 -- January 2026
\begin{itemize}
    \item Orchestrated 5 specialized agents on Google ADK with loop-guarded retry, per-agent model routing, and automated chart generation across 6 visualization types from natural language queries.
    \item Enforced read-only execution via 18 blocked SQL patterns and table allowlisting; abstracted dialect-specific rules across MySQL, PostgreSQL, and SQLite for cross-database portability.
\end{itemize}

\ressection{SKILLS \& ACHIEVEMENTS}

\textbf{Languages:} C++, Python, Java, JavaScript/TypeScript, SQL, Assembly, OCaml \\
\textbf{Backend \& Infrastructure:} FastAPI, Spring Boot, Kafka, Redis, Celery, Docker, AWS, Ray, Git \\
\textbf{ML \& Data:} PyTorch, CUDA, vLLM, HuggingFace Transformers, spaCy, Pandas, LangGraph \\
\textbf{Achievements:} Team China, International Science and Engineering Fair (ISEF); Silver Medal, Team China, International Young Physicists' Tournament (IYPT); Second Prize, China Mathematics Olympiad; Second Prize, China Physics Olympiad; Champion, Start-up in a Weekend Hackathon, Photon Hybrid Intelligence Track

\end{document}
